\section{Preliminaries}

\paragraph{Objective and clipping.}
We analyze optimization at a fixed reference dataset $D_0$, with public
size $n=|D_0|>0$, in dimension $d\geq1$. Record $a\in D_0$ has a differentiable
loss $f_a:\R^d\to\R$ and gradient $g_a=\nabla f_a$. For $C>0$, define
\[
 f(x)=\frac1n\sum_{a\in D_0}f_a(x),\qquad
 \bar g_C(x)=\frac1n\sum_{a\in D_0}\clip_C(g_a(x)),
 \qquad
 \clip_C(z)=
 \begin{cases}
 z,&\norm z\leq C,\\
 Cz/\norm z,&\norm z>C.
 \end{cases}
\]
The dataset size $n$ is distinct from the run-family index $j$.
Clipping always uses the Euclidean norm.

\paragraph{Private gradients.}
Fix a finite public record universe $\mathcal U\supseteq D_0$. For every
$a\in\mathcal U$, fix a query $g_a:\R^d\to\R^d$, agreeing with
$\nabla f_a$ on $D_0$ and common across private inputs.
Fix a Poisson sampling probability $0<q\leq1$ and a public denominator
$B>0$. In run $j$, the private gradient on an input $D\subseteq\mathcal U$ is
\begin{equation}\label{eq:private-gradient}
 g^{\mathrm{priv}}_{j,k}
 =\frac1B\sum_{a\in D}I_{j,k,a}
       \clip_C(g_a(x_{j,k-1}))
       +\frac{C\sigma_{\mathrm{DP},j}}B Z_{j,k},
 \qquad k=1,\ldots,T_j.
\end{equation}
The Gaussian noise multiplier satisfies $\sigma_{\mathrm{DP},j}>0$.
The masks are independent Bernoulli variables with parameter $q$, and the
$Z_{j,k}$ are independent $N(0,I_d)$ vectors, independent of all masks.
Each step uses fresh randomness; its query is evaluated at the parameter
determined by the preceding releases. The denominator is fixed across
privacy neighbors, not the realized batch size or a value recomputed from
the supplied dataset size.

Privacy uses add/remove adjacency on subsets of a finite public record
universe $\mathcal U\supseteq D_0$. Queries are fixed across inputs, and
initialization and algorithm parameters are public and common to
neighbors. Measurability of the release kernels is included in
Assumption~\ref{ass:queries}. The protected output is the
entire interpolated parameter path; masks and Gaussian seeds are not
released. Appendix~\ref{app:privacy} gives the precise mechanism setup,
reconstruction, and accounting argument.

\paragraph{DP-Adam.}
Fix a stabilizer $a_{\mathrm{Adam}}>0$ and momentum parameters
\begin{equation}\label{eq:betas}
 0<\beta_1<1,\qquad \beta_1^2<\beta_2<1.
\end{equation}
Starting both memories at zero, standard bias-corrected Adam in ordinary
gradient units is
\begin{align}
 m^{\mathrm{priv}}_{j,k}
 &=\beta_1m^{\mathrm{priv}}_{j,k-1}
        +(1-\beta_1)g^{\mathrm{priv}}_{j,k},\nonumber\\
 v^{\mathrm{priv}}_{j,k}
 &=\beta_2v^{\mathrm{priv}}_{j,k-1}
        +(1-\beta_2)(g^{\mathrm{priv}}_{j,k})^{\odot2},\label{eq:adam}\\
 x_{j,k}
 &=x_{j,k-1}-\eta_j
 \frac{m^{\mathrm{priv}}_{j,k}/(1-\beta_1^k)}
 {\sqrt{v^{\mathrm{priv}}_{j,k}/(1-\beta_2^k)}+a_{\mathrm{Adam}}}.
 \nonumber
\end{align}
Squares, square roots, and quotients are coordinatewise. The learning
rate $\eta_j$ is constant within a run; its dependence on $j$ compares
complete runs of different lengths, not a decay schedule within one run.

The proofs use the equivalent normalized variables
\[
 s_j=\frac{\sigma_{\mathrm{DP},j}}B,\qquad
 a_A=\frac{a_{\mathrm{Adam}}}C,\qquad
 m_{j,k}=\frac{m^{\mathrm{priv}}_{j,k}}C,\qquad
 v_{j,k}=\frac{v^{\mathrm{priv}}_{j,k}}{C^2}.
\]
Their input and the direction used by the algorithm are
\begin{align}
 Y_{j,k}:=\frac{g^{\mathrm{priv}}_{j,k}}C
 &=\frac1{CB}\sum_{a\in D}I_{j,k,a}
       \clip_C(g_a(x_{j,k-1}))+s_jZ_{j,k},\label{eq:release}\\
 \widehat A_{j,k}
 &=\frac{m_{j,k}/(1-\beta_1^k)}
 {\sqrt{v_{j,k}/(1-\beta_2^k)}+a_A},
 \qquad x_{j,k}=x_{j,k-1}-\eta_j\widehat A_{j,k}.
 \label{eq:normalized-direction}
\end{align}
The normalized memories obey the same recurrences with input $Y_{j,k}$.
The condition $\beta_1^2<\beta_2$ makes the weighted Cauchy--Schwarz
constant finite: in each coordinate the first moment is bounded by a
fixed multiple of the square root of the second moment. Consequently,
the normalized directions, including bias correction, are uniformly
bounded; see Lemma~\ref{lem:bounded-directions}.

\begin{assumption}[Time scaling and initialization]\label{ass:clock}
Let $T_j\in\mathbb N_{>0}$ and $\eta_j>0$. There are fixed $H,\lambda>0$
such that
\begin{equation}\label{eq:clock}
 T_j\eta_j=s_jH,\qquad s_j\longrightarrow\infty,
 \qquad \eta_js_j\longrightarrow\lambda.
\end{equation}
The initial points $x_j=x_{j,0}$ are public, common across privacy
neighbors, and satisfy $x_j\to x_0$.
\end{assumption}

Define the slow-time step $h_j=H/T_j=\eta_j/s_j$ and let $X_j$ affinely
interpolate the iterates on $[0,H]$:
\[
 X_j((k+u)h_j)=(1-u)x_{j,k}+u x_{j,k+1},
 \qquad 0\leq k<T_j,\quad 0\leq u\leq1.
\]
There are $T_j$ releases and $T_j+1$ nodes, including the terminal node.
The assumption implies $\eta_j\to0$, $h_j\to0$, and
\[
 \frac{s_j}{\sqrt{T_j}}\longrightarrow c_s:=\sqrt{\frac{\lambda}{H}},
 \qquad \eta_j\sqrt{T_j}\longrightarrow c_\eta:=\sqrt{H\lambda}.
\]
Consequently,
\[
 H=\frac{c_\eta}{c_s},\qquad\lambda=c_\eta c_s.
\]
The exact relation in \eqref{eq:clock} specifies the interpolation used
in the proof. The horizon $H$ is fixed before $j\to\infty$. Per-release
noise diverges, but the number of releases grows quadratically with it;
this does not imply vanishing overall privacy loss.

Throughout, $\norm{\cdot}$ is the Euclidean norm and $\snorm{\cdot}$
the coordinate maximum norm. For a path $w$, write
$\snorm w=\sup_{0\leq t\leq H}\snorm{w(t)}$.

\section{Theoretical results}

The main results identify the limiting dynamics, transfer a stationarity
bound to finite iterates, and use privacy calibration to select the
learning-rate coefficient that minimizes that bound. The dependence and
approximation arguments appear in the appendices.

\subsection{Limiting dynamics}

\paragraph{The two memory coefficients.}
It is enough to define them in one coordinate. Let $(Z_k)_{k\in\mathbb Z}$
be independent standard Gaussians, and set
\begin{equation}\label{eq:reference}
 \begin{aligned}
 p_\ell&=(1-\beta_1)\beta_1^\ell,
 &w_\ell&=(1-\beta_2)\beta_2^\ell,\\
 M_k&=\sum_{\ell\geq0}p_\ell Z_{k-\ell},
 &V_k&=\sum_{\ell\geq0}w_\ell Z_{k-\ell}^2,
 &U_k^G&=\frac{M_k}{\sqrt{V_k}}.
 \end{aligned}
\end{equation}
This stationary direction has mean zero. Let $U_k^G(r)$ denote the same
direction when every Gaussian input in its history is shifted by $r$.
Define the response coefficient and long-run variance by
\begin{equation}\label{eq:memory-coefficients}
 \alpha_\beta=
 \left.\frac{d}{dr}\E U_0^G(r)\right|_{r=0},\qquad
 \sigma_\beta^2=
 \operatorname{Var}(U_0^G)+2\sum_{\ell\geq1}
             \operatorname{Cov}(U_0^G,U_\ell^G).
\end{equation}
Thus $\alpha_\beta$ measures the response to a small common gradient
signal, including the change in the denominator. The coefficient
$\sigma_\beta^2$ includes both the variance of one update and the
covariance accumulated across iterations. Both are finite and strictly
positive. Independent, identically distributed coordinates give isotropic
long-run covariance $\sigma_\beta^2I_d$.
Appendix~\ref{app:reference} proves these facts and gives an explicit
expectation formula for $\alpha_\beta$.

\begin{assumption}[Fixed record queries]\label{ass:queries}
The public denominator satisfies $B=qn$, the expected minibatch size at
$D_0$, and is fixed across neighbors. It is not the realized Poisson batch
size. For each $a\in D_0$,
$g_a=\nabla f_a$, where
$f_a\in C^1(\R^d)$ and $g_a$ is globally Euclidean Lipschitz with constant
$L_a$. For every private input $D\subseteq\mathcal U$ and every
$j,k$, the release \eqref{eq:private-gradient} defines a measurable
probability kernel in the current state. No Lipschitz assumption is imposed
on queries for records outside $D_0$.
\end{assumption}
Put $L_C=n^{-1}\sum_{a\in D_0}L_a$. Nonexpansiveness of clipping gives
\[
 \|\bar g_C(x)\|_2\leq C,\qquad
 \|\bar g_C(x)-\bar g_C(y)\|_2\leq L_C\|x-y\|_2.
\]
Thus the limiting drift $-(\alpha_\beta/C)\bar g_C$ is bounded and globally
Lipschitz, independently of the objective bounds used below.

\begin{theorem}[High-noise diffusion limit]\label{thm:diffusion}
For the recursion \eqref{eq:adam}, suppose Assumptions~\ref{ass:clock}
and~\ref{ass:queries} hold. Then the laws of $X_j$ at $D_0$ converge weakly in
$C([0,H],\R^d)$ to the law of
\begin{equation}\label{eq:sde}
 dX_t=-\frac{\alpha_\beta}{C}\bar g_C(X_t)\,dt
       +\sqrt{\lambda}\,\sigma_\beta\,dW_t,\qquad X_0=x_0.
\end{equation}
Here $W$ is standard $d$-dimensional Brownian motion.
The drift is globally Lipschitz, so the additive equation has a unique
solution for each continuous driving path.
\end{theorem}

The limit separates the response and variance effects of Adam's memory.
The coefficient $\alpha_\beta$ controls how a small gradient signal changes
the mean update, while $\sigma_\beta^2$ includes the covariance accumulated across
correlated updates. Dominant Gaussian noise removes state-dependent
coordinatewise preconditioning at leading order. Replacing $\sigma_\beta^2$ by a
single-update variance would omit those temporal correlations. The
momentum pair remains fixed throughout the limit. No uniformity is
claimed as it approaches the boundary of the admissible region.

The second-moment estimate nevertheless remains random at the iteration
scale. In one coordinate of the Gaussian reference, independence and
$\operatorname{Var}(Z_{\ell}^2)=2$ give
\[
 \E V_0=1,\qquad
 \operatorname{Var}(V_0)=2\sum_{\ell\geq0}w_{\ell}^2
 =\frac{2(1-\beta_2)}{1+\beta_2}>0.
\]
The limiting coefficients account for these fluctuations and their temporal
dependence; the denominator is not replaced by its expectation.

The scaling follows from the frozen stationary mean direction $\bar U_j(x)$.
Lemma~\ref{lem:averaged-update} gives, uniformly in $x$ in the fixed-query setting,
\[
 \bar U_j(x)=\frac{\alpha_\beta}{Cs_j}\bar g_C(x)+O(s_j^{-2}).
\]
Since $h_j=\eta_j/s_j$, the drift and covariance scales are
\[
 \begin{aligned}
 -\frac{\eta_j\bar U_j(x)}{h_j}
 &=-s_j\bar U_j(x)\longrightarrow-\frac{\alpha_\beta}{C}\bar g_C(x),\\
 \frac{\eta_j^2}{h_j}\sigma_\beta^2I_d
 &=\eta_js_j\sigma_\beta^2I_d\longrightarrow\lambda\sigma_\beta^2I_d.
 \end{aligned}
\]
Appendix~\ref{app:transfer} connects this frozen calculation to the
adaptive trajectory.

\paragraph{Proof idea.}
The rescaled private release
$Y_{j,k}/s_j=g^{\mathrm{priv}}_{j,k}/(Cs_j)$ is a standard Gaussian
vector plus a uniformly bounded perturbation of order $s_j^{-1}$.
Expanding the stationary mean update identifies the averaged drift
$b_j=s_j\bar U_j$ and its response coefficient $\alpha_\beta$.

To identify the accumulated noise, we assign each Gaussian input all of
its future contributions to the reference updates. Their partial sums
differ from a bounded-increment martingale by a uniformly bounded boundary
term. A functional central limit theorem then gives the full long-run
covariance $\sigma_\beta^2I_d$.

The remaining step connects these stationary calculations to initialized,
adaptive Adam. Over a short block we would like to hold the parameter fixed
and use the corresponding stationary mean. A stationary copy, however,
starts with a different history. Bounding that history mismatch separately
on every block requires long blocks, while controlling parameter motion
requires short ones. We instead compare each process with a Gaussian
reference driven by the same noise. The signal and rescaled stabilizer
change its direction by only $O(s_j^{-1})$ in mean, apart from the initial
transient. Comparing these small corrections retains geometric forgetting
of the histories and permits both errors to vanish. Appendix~\ref{app:transfer}
explains the compatible block choice $L_{\mathrm{blk}}\asymp
\sqrt{s_j\log s_j}$. Summable bias corrections complete the decomposition
\[
 X_j(t)=x_j+W_j(t)-\int_0^t b_j(X_j(u))\,du+R_j(t),
 \qquad \E\snorm{R_j}\longrightarrow0.
\]
The reference functional limit, local uniform convergence
$b_j\to(\alpha_\beta/C)\bar g_C$, and continuity of the additive integral
equation identify the limiting diffusion.

\paragraph{Time measured as a fraction of the run.}
The same limit has a useful form on $[0,1]$. Set $\widetilde X_u=X_{Hu}$
and $\widetilde W_u=H^{-1/2}W_{Hu}$. Brownian time scaling gives
\begin{equation}\label{eq:run-time-sde}
 d\widetilde X_u
 =-\frac{\alpha_\beta c_\eta}{Cc_s}\bar g_C(\widetilde X_u)\,du
       +c_\eta\sigma_\beta\,d\widetilde W_u,
 \qquad \widetilde X_0=x_0.
\end{equation}
Here $\widetilde W$ is standard Brownian motion on $[0,1]$. At fixed
$c_\eta$, increasing $c_s$ reduces the coefficient of the gradient-driven
drift, while the Brownian coefficient remains unchanged.

\subsection{Stationarity and finite-iterate transfer}

Clipping bounds each record contribution but can change the direction of
the average gradient. We quantify that discrepancy through per-example gradient
disagreement. Nonexpansiveness of clipping then gives a descent inequality
with three contributions: the initial objective gap, per-example gradient disagreement,
and accumulated update noise.

\begin{assumption}[Objective and per-example gradient disagreement]\label{ass:objective}
The empirical objective $f\in C^2(\R^d)$ satisfies $f\geq\ell$.
For a constant $0\leq\Delta\leq C/2$,
\begin{equation}\label{eq:dispersion}
 \frac1{n}\sum_{a\in D_0}\norm{g_a(x)-\nabla f(x)}
 \leq\Delta\leq C/2\qquad\text{for all }x,
\end{equation}
and assume the one-sided Hessian bound
\begin{equation}\label{eq:hessian}
 D^2f(x)[z,z]\leq L_f\norm z^2
 \qquad\text{for all }x,z,\qquad L_f\geq0.
\end{equation}
\end{assumption}
In this fixed-query setting, $\nabla f$ is already globally Lipschitz,
with constant at most $L_C$, because it averages the record gradients.

The disagreement assumption is global and can be restrictive: for quadratic
record losses on $\R^d$, all record Hessians must coincide. Equal-Hessian
shifted quadratics can still have nonzero disagreement; Appendix~\ref{app:utility}
explains the restriction.

No convexity is required. Define the clipping-adapted stationarity measure
\begin{equation}\label{eq:penalty}
 P_C(x)=\min\{\norm{\nabla f(x)}^2,C\norm{\nabla f(x)}\}.
\end{equation}
It is quadratic when the gradient is below the clipping threshold and
linear beyond it. Small values indicate approximate stationarity, but the
bound gives weaker control of large gradients than a bound on the full
squared gradient norm. Its role in descent follows from the radial identity
\[
 P_C(x)=\langle\nabla f(x),\clip_C(\nabla f(x))\rangle.
\]
Its relation to approximate stationarity is explicit: for any random
vector $Z$ with $\E P_C(Z)<\infty$ and any $a>0$,
\begin{equation}\label{eq:stationarity-probability}
 \Prob(\norm{\nabla f(Z)}\geq a)
 \leq\frac{\E P_C(Z)}{\min\{a^2,Ca\}}.
\end{equation}
This follows from the monotonicity of $r\mapsto\min\{r^2,Cr\}$ and
Markov's inequality.

The algorithm averages clipped record gradients. Nonexpansiveness of
clipping and the disagreement assumption connect this average to the
clipped full gradient:
\[
 \|\bar g_C(x)-\clip_C(\nabla f(x))\|_2\leq\Delta.
\]
Writing $r=\norm{\nabla f(x)}$, this gives
$\langle\nabla f(x),\bar g_C(x)\rangle\geq\min\{r^2,Cr\}-\Delta r$.
For $r\leq C$, completing the square bounds the disagreement term; for
$r>C$, the condition $\Delta\leq C/2$ preserves half the clipped descent.
Consequently,
\[
 \boxed{\langle\nabla f(x),\bar g_C(x)\rangle
       \geq\tfrac12P_C(x)-\tfrac12\Delta^2.}
\]
Lemma~\ref{lem:alignment} gives the two-case calculation.

\begin{theorem}[Asymptotic stationarity bound for DP-Adam iterates]\label{cor:finite}
Under Assumptions~\ref{ass:clock}--\ref{ass:objective}, let $K_j$ be
independent of training and uniform on $\{0,\ldots,T_j-1\}$. At $D_0$,
\begin{equation}\label{eq:finite-rate}
\begin{aligned}
 \lim_{j\to\infty}\E P_C(x_{j,K_j})
 &=\frac1H\E\int_0^H P_C(X_t)\,dt\\
 &\leq
 \underbrace{\frac{2C(f(x_0)-\ell)}{\alpha_\beta H}}
 _{\text{initial objective gap}}
 +\underbrace{\Delta^2}_{\text{per-example gradient disagreement}}\\
 &+\underbrace{\frac{\lambda L_fCd\,\sigma_\beta^2}{\alpha_\beta}}
 _{\text{accumulated update noise}}.
\end{aligned}
\end{equation}
The expected stationarity measure is finite. Appendix~\ref{app:finite} proves the
occupation and grid-average expectation limits in this statement.
\end{theorem}

\paragraph{Descent calculation.}
For the generator $\mathcal L$ of the limiting diffusion, clipping
alignment and the upper Hessian bound give
\[
 \mathcal Lf
 \leq-\frac{\alpha_\beta}{2C}P_C
       +\frac{\alpha_\beta\Delta^2}{2C}
       +\frac{\lambda L_fd \sigma_\beta^2}{2}.
\]
Stopped It\^o calculus integrates this inequality; lower boundedness of
$f$ and Fatou's lemma then bound the diffusion occupation average.
Appendix~\ref{app:utility} supplies the localization and integrability
details and the diffusion-only bound. The moment and grid-average
arguments in Appendix~\ref{app:finite} then give the expectation equality above.

The bound separates three contributions. The initial objective gap decays
as $H^{-1}$. Per-example gradient disagreement contributes $\Delta^2$, through the
possible misalignment introduced by clipping. Accumulated noise contributes
$\lambda L_fCd\,\sigma_\beta^2/\alpha_\beta$, so Adam's memory affects both
variance and response to the gradient signal.

The proof in Appendix~\ref{app:finite} uses the
moment estimates behind the diffusion limit; weak convergence alone would
not suffice. Corollary~\ref{cor:objective-expectation} in Appendix~\ref{app:moments}
also gives convergence of expected terminal objective values.

\paragraph{A fixed-horizon clipped quadratic example.}
To illustrate the limiting dynamics at $H=1$, take one record with query
$g(x)=x$, objective $f(x)=x^2/2$, and initial
parameter $x_0=2$. Set $q=B=C=a_A=H=1$ and
$\beta_1=\beta_2=1/2$. The query is globally Lipschitz and the objective assumptions hold with
$\Delta=0$, $L_f=1$, and $\ell=0$. Choose any fixed $\lambda>0$.
Appendix~\ref{app:clock-examples} supplies a privacy-calibrated choice
and checks its schedule and release mechanism. The limiting equation
and its rate are
\[
 dX_t=-\alpha_\beta\clip_1(X_t)\,dt
             +\sqrt{\lambda}\,\sigma_\beta\,dW_t,\qquad X_0=2,
\]
\begin{equation}\label{eq:example}
 \frac1t\E\int_0^t\min\{X_u^2,|X_u|\}\,du
       \leq\frac{4}{\alpha_\beta t}
                     +\frac{\lambda\sigma_\beta^2}{\alpha_\beta},
 \qquad 0<t\leq1.
\end{equation}
The restoring drift is linear inside $[-1,1]$ and has constant magnitude
outside that interval. Clipping therefore remains nonlinear even though
Adam's coordinatewise normalization has averaged out of the parameter
equation. The long-run variance determines the noise intensity.

\subsection{Privacy calibration and learning-rate tuning}

Privacy accounting uses the same finite algorithms. Fix an integer order
$r\geq2$ and $0<\delta<1$. Let $\rho_r(q,\sigma)$ be the
sufficient sampled-Gaussian R\'enyi bound defined explicitly in
\eqref{eq:rdp}, and put
\[
 c_\delta=\frac{\log(1/\delta)}{r-1}.
\]

Proposition~\ref{prop:privacy} gives the finite accountant
$\varepsilon_j=T_j\rho_r(q,s_jB)+c_\delta$. Since
$c_s^2=\lambda/H$ and $B=qn$, its limit is
\[
 \varepsilon_j\longrightarrow\frac{r}{2n^2c_s^2}+c_\delta.
\]
Thus the limiting sufficient privacy bound depends on the noise
coefficient $c_s$. For a target $\varepsilon>c_\delta$ and public slack
$\zeta>0$, choose
\[
 c_s=\sqrt{\gamma_\varepsilon},\qquad
 \gamma_\varepsilon=\frac{(1+\zeta)r}{2n^2(\varepsilon-c_\delta)}.
\]
This calibration makes the limit strictly smaller than $\varepsilon$, so
all sufficiently large members satisfy the target privacy guarantee.

\paragraph{An explicit finite-run schedule.}
For any sequence of positive integers $T_j\to\infty$ and any $c_\eta>0$, set
\begin{equation}\label{eq:explicit-schedule}
 \boxed{
 s_j=\sqrt{\gamma_\varepsilon T_j},\qquad
 \eta_j=\frac{c_\eta}{\sqrt{T_j}},\qquad
 \sigma_{\mathrm{DP},j}=B\sqrt{\gamma_\varepsilon T_j}.}
\end{equation}
These choices satisfy the time-scaling relations in
Assumption~\ref{ass:clock} exactly, with
\[
 H=\frac{c_\eta}{\sqrt{\gamma_\varepsilon}},\qquad
 \lambda=c_\eta\sqrt{\gamma_\varepsilon}.
\]
This calibration sets the noise coefficient; $c_\eta$ is the remaining
learning-rate coefficient. More general families satisfying
Assumption~\ref{ass:clock} with these $H,\lambda$ are also covered.

\paragraph{Minimizing the descent bound.}
The nondegenerate diffusion regime imposes the $T_j^{-1/2}$ learning-rate
scaling; we minimize the bound over its coefficient $c_\eta$.
Fix public bounds $G,L>0$ such that $f(x_0)-\ell\leq G$ and $L_f\leq L$.
Substituting $H$ and $\lambda$ into Theorem~\ref{cor:finite} gives
\begin{equation}\label{eq:coefficient-bound}
 \boxed{\lim_{j\to\infty}\E P_C(x_{j,K_j})
 \leq\Delta^2+\frac{C\sqrt{\gamma_\varepsilon}}{\alpha_\beta}
 \left(\frac{2G}{c_\eta}+Ld\,\sigma_\beta^2 c_\eta\right).}
\end{equation}
The expectation limit exists by Theorem~\ref{cor:finite}.
The first term in parentheses decreases with $c_\eta$ and the second
increases. Privacy multiplies both by the same positive factor, so it
does not change the minimizing coefficient:
\[
 c_\eta^\star=\sqrt{\frac{2G}{Ld\,\sigma_\beta^2}},\qquad
 \Delta^2+\frac{2C\sigma_\beta}{\alpha_\beta}
          \sqrt{2GLd\,\gamma_\varepsilon}.
\]
The second expression is the minimized bound. After learning-rate tuning,
Adam's memory enters through the ratio $\sigma_\beta/\alpha_\beta$:
fluctuation relative to response. The corresponding diffusion
horizon is $H_*=c_\eta^\star/\sqrt{\gamma_\varepsilon}$.
The following theorem applies this choice to the finite algorithms.

\begin{theorem}[Learning-rate tuning under privacy calibration]\label{cor:horizon-tuned}
Fix $r\geq2$, $0<\delta<1$, $\varepsilon>c_\delta$, and
$\zeta>0$, with $r$ an integer.
Let $G,L>0$ be public bounds with $f(x_0)-\ell\leq G$ and $L_f\leq L$.
Using $\gamma_\varepsilon$ above, set
\[
 c_\eta^\star=\sqrt{\frac{2G}{Ld\,\sigma_\beta^2}},\qquad
 H_*=\frac{c_\eta^\star}{\sqrt{\gamma_\varepsilon}},\qquad
\lambda_*=c_\eta^\star\sqrt{\gamma_\varepsilon}.
\]
Suppose the algorithm family satisfies
Assumptions~\ref{ass:clock}--\ref{ass:objective} with $H=H_*$ and
$\lambda=\lambda_*$. Let $K_j$ be independent of training
and uniform on $\{0,\ldots,T_j-1\}$. For every tolerance $u>0$, all sufficiently
large members are $(\varepsilon,\delta)$-DP and satisfy
\[
 \E P_C(x_{j,K_j})
 \leq \Delta^2+\frac{2C\sigma_\beta}{\alpha_\beta}
      \sqrt{2GLd\,\gamma_\varepsilon}+u.
\]
The privacy threshold is data-independent; the utility threshold may
depend on $D_0$ and $u$. Every such family also satisfies
\[
 \eta_j\sqrt{T_j}\longrightarrow c_\eta^\star.
\]
In particular, \eqref{eq:explicit-schedule} with $c_\eta=c_\eta^\star$
has $\eta_j=\sqrt{2G/(Ld\,\sigma_\beta^2T_j)}$ exactly.
\end{theorem}
\begin{proof}
Under $B=qn$, the fixed-horizon calibration \eqref{eq:chosen} is
$\lambda_H=\gamma_\varepsilon H$. Thus
$\lambda_*=\lambda_H$ at $H=H_*$. Theorem~\ref{thm:main} gives eventual
privacy and the stationarity bound at this horizon. Replacing the initial
gap and Hessian coefficient by their public bounds $G,L$ gives
\eqref{eq:coefficient-bound} with $c_\eta=c_\eta^\star$, whose minimized
value is displayed above. The exact clock also gives
$\eta_j\sqrt{T_j}\to H_*\sqrt{\gamma_\varepsilon}=c_\eta^\star$ for every
family satisfying the stated assumptions.
\end{proof}

\paragraph{Interpretation and scope.}
The theorem selects the bound-minimizing learning-rate coefficient with
momentum, clipping, sampling, and public objective bounds fixed. The separate
growth estimate can be smaller, so this minimization need not optimize actual
stationarity or give a sharp privacy--utility rate; Appendix~\ref{sec:tuned-budget}
compares the two bounds.

Nor does a privacy-independent coefficient imply preserved descent.
For the small-budget calibration of that appendix,
$c_s=\sqrt{\gamma_\varepsilon}\to\infty$, while $c_\eta^\star$ stays fixed.
Realize the budget-dependent solutions $\widetilde X^\varepsilon$ of
\eqref{eq:run-time-sde} with the same Brownian driver $\widetilde W$.
The equation has vanishing drift and fixed noise intensity. Using
$\norm{\bar g_C}\leq C$ gives, almost surely,
\begin{equation}\label{eq:brownian-comparison}
 \sup_{0\leq u\leq1}
 \norm{\widetilde X^\varepsilon_u-
       \bigl(x_0+c_\eta^\star\sigma_\beta\widetilde W_u\bigr)}
 \leq\frac{\alpha_\beta c_\eta^\star}{c_s}.
\end{equation}
The algorithm limit is taken separately at each fixed privacy target before
this small-budget comparison. Utility concerns a uniformly sampled iterate
at $D_0$, while privacy holds across neighbors. The vanishing utility error
is unquantified: the adaptive-remainder estimate is not a rate for the full
path-law approximation.
